\documentclass[11pt]{article}

% Basic packages
\usepackage[margin=1in]{geometry}
\usepackage{amsmath, amssymb, amsthm, mathtools}
\usepackage{bm}
\usepackage{enumerate}
\usepackage{enumitem}
\usepackage{graphicx}
\usepackage{booktabs}
\usepackage{hyperref}
\hypersetup{
    colorlinks=true,
    linkcolor=blue,
    urlcolor=blue,
    citecolor=blue
}

\setlength{\parskip}{0.5em}
\setlength{\parindent}{0pt}

\begin{document}

\begin{center}
    {\LARGE \textbf{AMS 598 Final Exam Lecture Notes}}\\[0.3cm]
    {\large MapReduce, ADMM, Logistic Regression, SVM, Random Forests}\\[0.2cm]
\end{center}

\tableofcontents
\newpage

%%%%%%%%%%%%%%%%%%%%%%%%%%%%%%%%%%%%%%%%%%%%%%%%%%%%%%%%%%%%
\section{MapReduce Algorithms}
%%%%%%%%%%%%%%%%%%%%%%%%%%%%%%%%%%%%%%%%%%%%%%%%%%%%%%%%%%%%

\subsection{The MapReduce Paradigm}

MapReduce is a programming model for large-scale data processing on a cluster.

\begin{itemize}
    \item Input data is split into \emph{key--value pairs} $(k, v)$.
    \item \textbf{Map phase:} Apply a map function to each input pair, producing intermediate key--value pairs:
    \[
        (k, v) \mapsto \{(k_1, v_1), (k_2, v_2), \dots\}.
    \]
    \item The system groups all intermediate values by key.
    \item \textbf{Reduce phase:} For each key $k$, apply a reduce function to the list of values:
    \[
        (k, [v_1, v_2, \dots]) \mapsto (k, v_{\text{out}}).
    \]
    \item Tasks:
    \begin{itemize}
        \item Map tasks run in parallel on different data chunks.
        \item Reduce tasks run in parallel on different keys.
    \end{itemize}
\end{itemize}

Key idea: Make each map and reduce operation independent so that they can be parallelized.

%%%%%%%%%%%%%%%%%%%%%%%%%%%%%%%%%%%%%%%%%%%%%%%%%%%%%%%%%%%%
\subsection{TF--IDF in MapReduce}

TF--IDF: Term Frequency--Inverse Document Frequency.

\subsubsection*{Motivation}

We have a collection of $N$ documents. For a term $i$ and a document $j$:
\begin{itemize}
    \item $f_{ij}$: frequency (count) of term $i$ in document $j$.
    \item $n_i$: number of documents in which term $i$ appears.
\end{itemize}

Definitions:
\[
    \text{TF}_{ij} = \frac{f_{ij}}{\max_k f_{kj}},\qquad
    \text{IDF}_i = \log_2 \frac{N}{n_i},\qquad
    \text{TFIDF}_{ij} = \text{TF}_{ij} \cdot \text{IDF}_i.
\]

Goals:\,
\begin{itemize}
    \item Down-weight \emph{very frequent} terms (e.g.\ ``the'').
    \item Up-weight terms that are concentrated in relatively few documents and thus more indicative of topic.
\end{itemize}

\subsubsection*{Information Needed}

To compute TF--IDF, we need:
\begin{enumerate}[label=(\arabic*)]
    \item Number of times term $i$ appears in each document $j$ ($f_{ij}$).
    \item Maximum frequency in each document $j$ (i.e.\ $\max_k f_{kj}$).
    \item Number of documents containing term $i$ ($n_i$).
    \item Total number of documents $N$.
\end{enumerate}

This naturally leads to a multi-pass MapReduce pipeline.

\subsubsection*{Job 1: Word Frequency in Each Document}

\textbf{Input:} (docname, contents) pairs.

\textbf{Mapper:}
\begin{itemize}
    \item For each word $\texttt{w}$ in document $d$, emit:
    \[
        (d, \texttt{contents}) \mapsto ((\texttt{w}, d), 1).
    \]
\end{itemize}

\textbf{Reducer:}
\begin{itemize}
    \item For each key $(\texttt{w}, d)$, sum the counts:
    \[
        ((\texttt{w}, d), [1,1,\dots,1]) \mapsto ((\texttt{w}, d), f_{wd}).
    \]
\end{itemize}

Output: $((\texttt{w}, d), f_{wd})$.

\subsubsection*{Job 2: Max Word Count per Document}

\textbf{Input:} $((\texttt{w}, d), f_{wd})$ from Job 1.

\textbf{Mapper:}
\[
    ((\texttt{w}, d), f_{wd}) \mapsto (d, (\texttt{w}, f_{wd})).
\]

\textbf{Reducer:}
\begin{itemize}
    \item For each document $d$, we receive $(\texttt{w}, f_{wd})$ for many words.
    \item Compute $m_d = \max_{\texttt{w}} f_{wd}$.
    \item For each word in this doc, re-emit:
    \[
        ((\texttt{w}, d), (f_{wd}, m_d)).
    \]
\end{itemize}

Now we have per-document term counts and max frequency.

\subsubsection*{Job 3: Document Frequency per Term}

\textbf{Input:} $((\texttt{w}, d), (f_{wd}, m_d))$.

\textbf{Mapper:}
\[
    ((\texttt{w}, d), (f_{wd}, m_d)) \mapsto (\texttt{w}, (d, f_{wd}, m_d, 1)).
\]

\textbf{Reducer:}
\begin{itemize}
    \item For each word $\texttt{w}$, receive $(d, f_{wd}, m_d, 1)$ for all documents containing $\texttt{w}$.
    \item Sum all the $1$'s to get $n_{\texttt{w}}$, the number of documents containing $\texttt{w}$.
    \item Emit for each doc:
    \[
        ((\texttt{w}, d), (f_{wd}, m_d, n_{\texttt{w}})).
    \]
\end{itemize}

\subsubsection*{Job 4: Compute TF--IDF}

Assume $N$ is known.

\textbf{Input:} $((\texttt{w}, d), (f_{wd}, m_d, n_{\texttt{w}}))$.

\textbf{Mapper:}
\begin{align*}
    \text{TF}_{\texttt{w}d} &= \frac{f_{wd}}{m_d},\\
    \text{IDF}_{\texttt{w}} &= \log_2\frac{N}{n_{\texttt{w}}}.
\end{align*}
Emit:
\[
    ((\texttt{w}, d), \text{TF}_{\texttt{w}d} \cdot \text{IDF}_{\texttt{w}}).
\]

\textbf{Reducer:} Identity (just passes values through).

\subsubsection*{Scalability Remarks}

\begin{itemize}
    \item Very high-frequency words (e.g.\ ``the'') can cause memory issues when summing $n_i$.
    \item Possible solutions:
    \begin{itemize}
        \item Ignore terms exceeding a frequency threshold.
        \item Write intermediate results to disk or use more passes.
    \end{itemize}
    \item Jobs 3 and 4 can sometimes be combined to save I/O.
\end{itemize}

%%%%%%%%%%%%%%%%%%%%%%%%%%%%%%%%%%%%%%%%%%%%%%%%%%%%%%%%%%%%
\subsection{Relational Algebra Operations in MapReduce}

We can express standard relational operations (selection, projection, join, etc.) via MapReduce by choosing keys appropriately.

Let $R$ be a relation (table) with tuples $t$.

\subsubsection*{Selection $\sigma_C(R)$}

Condition $C(t)$ on each tuple.

\textbf{Map:}
\[
    t\in R \mapsto
    \begin{cases}
        (t, t), & \text{if } C(t)\text{ is true},\\
        \text{emit nothing}, & \text{otherwise}.
    \end{cases}
\]

\textbf{Reduce:} Identity.  
Result set: all selected tuples.

\subsubsection*{Projection $\pi_S(R)$}

We want only attributes in subset $S$, and we want duplicates removed.

\textbf{Map:}
\begin{itemize}
    \item For each tuple $t$, form $t'$ by projecting onto $S$.
    \item Emit $(t', t')$.
\end{itemize}

\textbf{Reduce:}
\[
    (t', [t', t', \dots, t']) \mapsto (t', t').
\]
Exactly one copy of each distinct $t'$ remains.

\subsubsection*{Union and Intersection}

Assume $R$ and $S$ have the same schema.

\paragraph{Union $R \cup S$:}
\textbf{Map:} For each tuple $t$ in $R$ or $S$, emit $(t,t)$.  
\textbf{Reduce:} Identity; we keep one copy of each $t$.

\paragraph{Intersection $R \cap S$:}
\textbf{Map:} Same as union.  
\textbf{Reduce:}
\[
    (t, [t,t]) \mapsto (t,t); \quad (t,[t]) \mapsto \text{nothing}.
\]
Only tuples that appear in both relations survive.

\subsubsection*{Difference $R - S$}

\textbf{Map:}
\begin{itemize}
    \item If $t \in R$, emit $(t, R)$.
    \item If $t \in S$, emit $(t, S)$.
\end{itemize}

\textbf{Reduce:}
\[
    (t, \text{values}) \mapsto
    \begin{cases}
        (t,t), & \text{if values} = [R],\\
        \text{nothing}, & \text{otherwise (} [R,S], [S], [S,R] \text{)}.
    \end{cases}
\]

\subsubsection*{Natural Join}

Example: $R(A,B)$ and $S(B,C)$, joining on attribute $B$.

\textbf{Map:}
\begin{itemize}
    \item For each $(a,b)\in R$, emit $(b, (R,a))$.
    \item For each $(b,c)\in S$, emit $(b, (S,c))$.
\end{itemize}

\textbf{Reduce:}
\begin{itemize}
    \item For each key $b$, collect two lists:
    \[
        \{(R,a_i)\},\quad \{(S,c_j)\}.
    \]
    \item Output all combinations $(a_i, b, c_j)$.
\end{itemize}

%%%%%%%%%%%%%%%%%%%%%%%%%%%%%%%%%%%%%%%%%%%%%%%%%%%%%%%%%%%%
\subsection{Matrix Multiplication in MapReduce}

Consider matrices $M$ and $N$:
\[
    M \in \mathbb{R}^{I\times J},\quad N \in \mathbb{R}^{J\times K},\quad P = MN \in \mathbb{R}^{I\times K}.
\]
Elements:
\[
    p_{ik} = \sum_{j} m_{ij} n_{jk}.
\]

We represent matrices as relations:
\[
    M(I,J,V) = \{(i,j, m_{ij})\},\quad N(J,K,W) = \{(j,k, n_{jk})\}.
\]

\subsubsection*{Two-Pass MapReduce Algorithm}

\paragraph{First pass: join on $j$ and compute partial products.}

\textbf{Map:}
\begin{itemize}
    \item For each $(i,j,v)$ in $M$, emit $(j, (M, i, v))$.
    \item For each $(j,k,w)$ in $N$, emit $(j, (N, k, w))$.
\end{itemize}

\textbf{Reduce (per $j$):}
\begin{itemize}
    \item For each $j$, obtain:
    \[
        \{(M, i, m_{ij})\},\quad \{(N, k, n_{jk})\}.
    \]
    \item For all $i$ and $k$ pairs, compute $m_{ij}n_{jk}$ and emit
    \[
        ((i,k), m_{ij}n_{jk}).
    \]
\end{itemize}

\paragraph{Second pass: sum partial products.}

\textbf{Map:} Identity on keys $(i,k)$:
\[
    ((i,k), v) \mapsto ((i,k), v).
\]

\textbf{Reduce:}
\[
    ((i,k), [v_1, v_2, \dots]) \mapsto ((i,k), \sum_q v_q) = ((i,k), p_{ik}).
\]

\subsubsection*{One-Pass Approach (Conceptual)}

Emit all necessary keys $(i,k)$ directly from the mapper, so that each reducer can match on $j$ and sum the products. This is more demanding on memory and network but can be done when dimensions are manageable.

%%%%%%%%%%%%%%%%%%%%%%%%%%%%%%%%%%%%%%%%%%%%%%%%%%%%%%%%%%%%
\subsection{PageRank}

PageRank measures the stationary distribution of a random walk on the web graph.

\subsubsection*{Basic PageRank Without Teleportation}

Let $M$ be the transition matrix of the web graph:
\begin{itemize}
    \item $M_{ij} = 1/\text{outdeg}(j)$ if there is a link from page $j$ to page $i$.
    \item Otherwise, $M_{ij} = 0$.
\end{itemize}

If $v^{(t)}$ is the probability distribution of the surfer at step $t$, then:
\[
    v^{(t+1)}) = M v^{(t)}.
\]

Under certain conditions (strong connectivity, no dead ends), $v^{(t)}$ converges to a stationary vector $v$ satisfying:
\[
    v = Mv.
\]

\subsubsection*{Problems: Dead Ends \& Spider Traps}

\begin{itemize}
    \item \textbf{Dead end:} A page with no outgoing links $\Rightarrow$ mass ``leaks''.
    \item \textbf{Spider trap:} A subset of pages that only link among themselves $\Rightarrow$ mass gets stuck.
\end{itemize}

\subsubsection*{Taxation / Teleportation}

Solution: With probability $\beta$ follow a link, with probability $1-\beta$ teleport to a random page.

Let $e$ be the vector of all ones, and $n$ the number of pages. The PageRank iteration becomes:
\[
    v^{(t+1)} = \beta M v^{(t)} + (1 - \beta)\frac{1}{n} e.
\]

This always converges to a unique stationary distribution.

\subsubsection*{MapReduce Implementation Outline}

\paragraph{Phase 1: Build graph representation}

Input: (URL, HTML).  
Mapper:
\begin{itemize}
    \item Parse outlinks (list of URLs).
    \item Initialize rank $PR_{\text{init}} = 1/n$.
    \item Emit $(\text{URL}, (PR_{\text{init}}, \text{outlink-list}))$.
\end{itemize}

Reducer: identity.

\paragraph{Phase 2: Iterative PageRank}

In each iteration:

\textbf{Mapper:} For each page:
\[
    (\text{URL}, (\text{PR}, \text{outlinks})) \mapsto
\begin{cases}
    (\text{u}, \text{PR}/|\text{outlinks}|) & \text{for each u in outlinks},\\
    (\text{URL}, \text{outlinks}) & \text{to carry adjacency}.
\end{cases}
\]

\textbf{Reducer:}
\begin{itemize}
    \item Sum rank contributions $s_{\text{URL}}$ from in-links.
    \item Apply taxation:
    \[
        \text{PR}_{\text{new}}(\text{URL}) = (1-\beta)\frac{1}{n} + \beta s_{\text{URL}}.
    \]
    \item Emit $(\text{URL}, (\text{PR}_{\text{new}}, \text{outlinks}))$.
\end{itemize}

Iterate until convergence (e.g., 50--75 iterations on the web).

%%%%%%%%%%%%%%%%%%%%%%%%%%%%%%%%%%%%%%%%%%%%%%%%%%%%%%%%%%%%
\section{Penalty-Based Methods and ADMM}
%%%%%%%%%%%%%%%%%%%%%%%%%%%%%%%%%%%%%%%%%%%%%%%%%%%%%%%%%%%%

We consider linear models with penalties to control complexity, and we use ADMM to solve large-scale versions.

%%%%%%%%%%%%%%%%%%%%%%%%%%%%%%%%%%%%%%%%%%%%%%%%%%%%%%%%%%%%
\subsection{Ridge Regression (L2 Penalty)}

Model: $y = X\beta + \varepsilon$, where $X \in \mathbb{R}^{N\times p}$, $y \in \mathbb{R}^{N}$.

Penalized objective:
\[
    \text{RSS}_\lambda(\beta) = \|y - X\beta\|_2^2 + \lambda \|\beta\|_2^2.
\]

Closed-form solution:
\[
    \hat{\beta}_{\text{ridge}} = (X^T X + \lambda I)^{-1} X^T y.
\]

Ridge shrinks coefficients towards zero but does not set them exactly to zero. It reduces variance and stabilizes estimates when predictors are correlated.

%%%%%%%%%%%%%%%%%%%%%%%%%%%%%%%%%%%%%%%%%%%%%%%%%%%%%%%%%%%%
\subsection{LASSO (L1 Penalty)}

Objective:
\[
    \min_\beta \ \|y - X\beta\|_2^2 + \lambda \|\beta\|_1.
\]

\begin{itemize}
    \item Encourages sparsity: many coefficients are exactly zero.
    \item Geometrically, the $\ell_1$ penalty has diamond-shaped level sets, which intersect quadratic loss contours at coordinate axes.
\end{itemize}

%%%%%%%%%%%%%%%%%%%%%%%%%%%%%%%%%%%%%%%%%%%%%%%%%%%%%%%%%%%%
\subsection{Penalized Logistic Regression}

Binary response $Y \in \{0,1\}$, features $x_i \in \mathbb{R}^p$.

Model:
\[
    \mathbb{P}(Y_i = 1 \mid x_i) = \pi_i = \sigma(x_i^T \beta),
\]
where $\sigma(z) = \frac{1}{1+e^{-z}}$ is the logistic sigmoid.

Log-likelihood:
\[
    \ell(\beta) = \sum_{i=1}^N \left[ y_i \log \pi_i + (1-y_i)\log(1-\pi_i) \right].
\]

Penalized objective (to be \emph{maximized} or equivalently \emph{minimize} minus log-likelihood):
\[
    \max_\beta \; \ell(\beta) - r(\beta)
    \quad \Leftrightarrow \quad
    \min_\beta \; -\ell(\beta) + r(\beta),
\]
where $r(\beta)$ can be $\lambda \|\beta\|_2^2$, $\lambda \|\beta\|_1$, or a custom penalty.

%%%%%%%%%%%%%%%%%%%%%%%%%%%%%%%%%%%%%%%%%%%%%%%%%%%%%%%%%%%%
\subsection{ADMM: Alternating Direction Method of Multipliers}

We want to solve a constrained convex problem of the form:
\[
    \min_{x,z} f(x) + g(z) \quad \text{s.t. } Ax + Bz = c.
\]

Augmented Lagrangian:
\[
    L_\rho(x, z, y) = f(x) + g(z) + y^T(Ax + Bz - c) + \frac{\rho}{2}\|Ax + Bz - c\|_2^2,
\]
where $y$ is the dual variable and $\rho > 0$ is a penalty parameter.

\subsubsection*{ADMM Iterations}

\begin{enumerate}
    \item \textbf{x-update:}
    \[
        x^{k+1} = \arg\min_x L_\rho(x, z^k, y^k).
    \]
    \item \textbf{z-update:}
    \[
        z^{k+1} = \arg\min_z L_\rho(x^{k+1}, z, y^k).
    \]
    \item \textbf{Dual update:}
    \[
        y^{k+1} = y^k + \rho (A x^{k+1} + B z^{k+1} - c).
    \]
\end{enumerate}

Scaled form: use $u = \frac{1}{\rho}y$ and rewrite updates to improve numerical stability.

%%%%%%%%%%%%%%%%%%%%%%%%%%%%%%%%%%%%%%%%%%%%%%%%%%%%%%%%%%%%
\subsection{Consensus Optimization via ADMM}

We want to minimize a sum of functions:
\[
    \min_\beta \ \sum_{i=1}^N f_i(\beta),
\]
where $f_i$ is the loss for partition $i$ of the data.

Introduce local copies $\beta_i$ and a global consensus variable $\beta$:
\[
    \min_{\beta_1,\dots,\beta_N,\beta} \ \sum_{i=1}^N f_i(\beta_i) + r(\beta)\quad \text{s.t. } \beta_i = \beta \ \forall i.
\]

This is in ADMM form:
\[
    f(\{\beta_i\}) = \sum_i f_i(\beta_i),\quad g(\beta) = r(\beta),
\]
with constraints $\beta_i - \beta = 0$.

\subsubsection*{ADMM Updates for Consensus}

Let $u_i$ be scaled dual variables.

\paragraph{1. Local update (per partition):}
\[
    \beta_i^{k+1} = \arg\min_{\beta_i} \left[ f_i(\beta_i) + \frac{\rho}{2}\|\beta_i - \beta^k + u_i^k\|_2^2 \right].
\]
This is typically a penalized regression or logistic regression using only data from partition $i$, shrunk toward $(\beta^k - u_i^k)$.

\paragraph{2. Consensus update:}
Define
\[
    \bar{\beta}^{k+1} = \frac{1}{N} \sum_{i=1}^N \left(\beta_i^{k+1} + u_i^k\right).
\]
Then
\[
    \beta^{k+1} = \arg\min_\beta \left[ r(\beta) + \frac{N \rho}{2} \|\beta - \bar{\beta}^{k+1}\|_2^2 \right].
\]

\begin{itemize}
    \item If $r(\beta) = 0$, then $\beta^{k+1} = \bar{\beta}^{k+1}$.
    \item If $r(\beta) = \lambda \|\beta\|_2^2$ (ridge), then this is a quadratic minimization with closed form.
    \item If $r(\beta) = \lambda \|\beta\|_1$ (lasso), then this is a soft-thresholding step applied to $\bar{\beta}^{k+1}$.
\end{itemize}

\paragraph{3. Dual update:}
\[
    u_i^{k+1} = u_i^k + \left(\beta_i^{k+1} - \beta^{k+1}\right).
\]

\subsubsection*{Large-Scale Logistic Regression via ADMM}

\begin{itemize}
    \item Split data into $K$ partitions: $(X_i, y_i)$.
    \item Local objective on partition $i$:
    \[
        f_i(\beta_i) = -\ell_i(\beta_i) = -\sum_{j \in \text{part } i} \left[y_j \log \pi_j + (1-y_j)\log(1-\pi_j)\right],
    \]
    where $\pi_j = \sigma(x_j^T \beta_i)$.
    \item ADMM coordinates local logistic regressions through consensus $\beta$.
\end{itemize}

This allows logistic regression on massive data using parallel computing.

%%%%%%%%%%%%%%%%%%%%%%%%%%%%%%%%%%%%%%%%%%%%%%%%%%%%%%%%%%%%
\subsection{ADMM for Logistic Regression with a Custom Penalty Term}

The exam may explicitly ask: \emph{``Write out the ADMM algorithm for logistic regression with an L1, L2, or custom penalty term.''} This section gives the general template plus common special cases.

\subsubsection*{Problem Setup}

We have data split across $K$ partitions:
\[
    \{(X_i, y_i)\}_{i=1}^K.
\]
On partition $i$:
\[
    \ell_i(\beta_i) = -\sum_{j \in \text{part } i} \left[
        y_j \log \pi_j + (1 - y_j) \log (1 - \pi_j)
    \right], \quad \pi_j = \sigma(x_j^T \beta_i).
\]

We want to solve the penalized logistic regression:
\[
    \min_\beta \ \sum_{i=1}^K \ell_i(\beta) + r(\beta),
\]
where $r(\beta)$ is a \emph{custom penalty} term (can be L1, L2, elastic net, group penalty, etc).

\subsubsection*{Consensus Formulation}

Introduce local copies $\beta_i$ and a global variable $\beta$:
\[
    \min_{\{\beta_i\},\beta} \ \sum_{i=1}^K \ell_i(\beta_i) + r(\beta)
    \quad \text{s.t. } \beta_i = \beta, \; i = 1,\dots,K.
\]

This is exactly in the consensus ADMM form.

\subsubsection*{Augmented Lagrangian (Scaled Form)}

Introduce scaled dual variables $u_i$ and penalty parameter $\rho > 0$. The scaled-form augmented Lagrangian is:
\[
    L_\rho(\{\beta_i\}, \beta, \{u_i\}) =
    \sum_{i=1}^K \left[
        \ell_i(\beta_i)
        + \frac{\rho}{2} \|\beta_i - \beta + u_i\|_2^2
    \right]
    + r(\beta) + \text{const}.
\]

ADMM alternates between updating local $\beta_i$, global $\beta$, and dual $u_i$.

\subsubsection*{ADMM Updates (General Custom Penalty)}

\paragraph{1. Local logistic regression update (per partition $i$):}
\[
    \beta_i^{k+1}
    = \arg\min_{\beta_i}
    \left[
        \ell_i(\beta_i)
        + \frac{\rho}{2} \|\beta_i - \beta^k + u_i^k\|_2^2
    \right].
\]

Interpretation:
\begin{itemize}
    \item This is a logistic regression problem on partition $i$, with an additional quadratic penalty that shrinks $\beta_i$ toward $(\beta^k - u_i^k)$.
    \item In practice, you would solve this with a standard logistic regression solver (e.g.\ Newton, coordinate descent, or a library like \texttt{glmnet}) with an L2 penalty centered at $(\beta^k - u_i^k)$.
\end{itemize}

\paragraph{2. Consensus update for global $\beta$ (custom penalty):}

Compute the average of adjusted local variables:
\[
    \bar{\beta}^{k+1}
    = \frac{1}{K} \sum_{i=1}^K \left(\beta_i^{k+1} + u_i^k\right).
\]

Then solve:
\[
    \beta^{k+1}
    = \arg\min_\beta \left[
        r(\beta) + \frac{K \rho}{2} \|\beta - \bar{\beta}^{k+1}\|_2^2
    \right].
\]

This is the \textbf{proximal step} of the custom penalty:
\[
    \beta^{k+1} = \text{prox}_{\frac{1}{K\rho} r}(\bar{\beta}^{k+1}).
\]

\paragraph{3. Dual variable update:}
\[
    u_i^{k+1}
    = u_i^k + \left(\beta_i^{k+1} - \beta^{k+1}\right),
    \quad i = 1,\dots,K.
\]

These three lines are what you should be ready to write in the exam for the general custom-penalty case.

\subsubsection*{Special Cases: L2, L1, and Other Penalties}

\paragraph{(a) L2 penalty (ridge logistic regression):}
\[
    r(\beta) = \lambda \|\beta\|_2^2.
\]

Consensus update is:
\[
    \beta^{k+1} = \arg\min_\beta
    \left[
      \lambda \|\beta\|_2^2 + \frac{K\rho}{2} \|\beta - \bar{\beta}^{k+1}\|_2^2
    \right].
\]

This is a quadratic problem with closed form. Setting the gradient to zero:
\[
    2\lambda \beta + K\rho (\beta - \bar{\beta}^{k+1}) = 0
    \quad \Rightarrow \quad
    (2\lambda + K\rho)\beta = K\rho \,\bar{\beta}^{k+1}.
\]

So:
\[
    \beta^{k+1}
    = \frac{K\rho}{2\lambda + K\rho} \,\bar{\beta}^{k+1}.
\]

\paragraph{(b) L1 penalty (lasso logistic regression):}
\[
    r(\beta) = \lambda \|\beta\|_1.
\]

Consensus update is:
\[
    \beta^{k+1}
    = \arg\min_\beta
    \left[
        \lambda \|\beta\|_1 + \frac{K\rho}{2} \|\beta - \bar{\beta}^{k+1}\|_2^2
    \right].
\]

This is elementwise \textbf{soft-thresholding}:
\[
    \beta_j^{k+1}
    = \mathcal{S}_{\lambda/(K\rho)}\left(\bar{\beta}^{k+1}_j\right),
\]
where
\[
    \mathcal{S}_\tau(z) = \text{sign}(z)\,\max(|z| - \tau, 0)
\]
is the soft-thresholding operator.

\paragraph{(c) Elastic net or other separable penalties:}

For a separable penalty
\[
    r(\beta) = \sum_{j} r_j(\beta_j),
\]
the consensus step decouples across coordinates:
\[
    \beta_j^{k+1}
    = \arg\min_{\beta_j}
    \left[
        r_j(\beta_j) + \frac{K\rho}{2} (\beta_j - \bar{\beta}^{k+1}_j)^2
    \right].
\]
Each coordinate update is a 1D proximal step.

\paragraph{(d) Group penalties (e.g.\ group lasso):}
If
\[
    r(\beta) = \sum_{g \in \mathcal{G}} \lambda_g \|\beta_g\|_2,
\]
where $\beta_g$ is a block of coefficients, then the consensus step:
\[
    \min_\beta \sum_{g} \lambda_g \|\beta_g\|_2 + \frac{K\rho}{2} \sum_g \|\beta_g - \bar{\beta}_g^{k+1}\|_2^2
\]
decouples over groups. Each group update is a \emph{block soft-thresholding} (a shrinkage in the direction of $\bar{\beta}_g^{k+1}$).

\subsubsection*{Summary: ``Template Answer'' for the Exam}

If asked to \emph{``write out the ADMM algorithm for logistic regression with a custom penalty term''}, you can write:

\begin{itemize}
    \item Problem:
    \[
        \min_{\{\beta_i\},\beta}
        \sum_{i=1}^K \ell_i(\beta_i) + r(\beta)
        \quad \text{s.t. } \beta_i = \beta.
    \]
    \item ADMM iterations:
    \begin{align*}
        \beta_i^{k+1}
        &= \arg\min_{\beta_i} \left[
            \ell_i(\beta_i)
            + \frac{\rho}{2}\|\beta_i - \beta^k + u_i^k\|_2^2
        \right], \\
        \bar{\beta}^{k+1} &= \frac{1}{K} \sum_{i=1}^K (\beta_i^{k+1} + u_i^k),\\
        \beta^{k+1}
        &= \arg\min_\beta \left[
            r(\beta) + \frac{K\rho}{2} \|\beta - \bar{\beta}^{k+1}\|_2^2
        \right], \\
        u_i^{k+1} &= u_i^k + \left(\beta_i^{k+1} - \beta^{k+1}\right).
    \end{align*}
\end{itemize}

Then briefly state the closed-form solution of the consensus step for L1, L2, or whichever specific penalty the question mentions.

%%%%%%%%%%%%%%%%%%%%%%%%%%%%%%%%%%%%%%%%%%%%%%%%%%%%%%%%%%%%
\section{Support Vector Machines (SVM)}
%%%%%%%%%%%%%%%%%%%%%%%%%%%%%%%%%%%%%%%%%%%%%%%%%%%%%%%%%%%%

SVM is a maximum-margin classifier in a (possibly high-dimensional) feature space.

%%%%%%%%%%%%%%%%%%%%%%%%%%%%%%%%%%%%%%%%%%%%%%%%%%%%%%%%%%%%
\subsection{Linear Separating Hyperplane}

Given labeled data $(x_i, y_i)$ with $y_i \in \{-1, +1\}$, a linear classifier is:
\[
    f(x) = \text{sign}(\beta^T x + b).
\]

Hyperplane: $\beta^T x + b = 0$.

%%%%%%%%%%%%%%%%%%%%%%%%%%%%%%%%%%%%%%%%%%%%%%%%%%%%%%%%%%%%
\subsection{Margin and Optimal Separator}

Distance from a point $x$ to the hyperplane is:
\[
    r = \frac{\beta^T x + b}{\|\beta\|_2}.
\]

For a correctly classified point:
\[
    y_i(\beta^T x_i + b) > 0.
\]

\textbf{Margin} $\rho$ is the width between the two class boundaries:
\[
    \beta^T x + b = 1,\quad \beta^T x + b = -1.
\]
Distance between these planes:
\[
    \rho = \frac{2}{\|\beta\|_2}.
\]

Max-margin classifier $\Rightarrow$ minimize $\|\beta\|_2$ (or equivalently $\frac{1}{2}\|\beta\|_2^2$).

%%%%%%%%%%%%%%%%%%%%%%%%%%%%%%%%%%%%%%%%%%%%%%%%%%%%%%%%%%%%
\subsection{Hard-Margin SVM}

Assume data are linearly separable and scaled so the margin boundaries correspond to $\pm1$.

Primal problem:
\[
    \min_{\beta,b} \ \frac{1}{2}\|\beta\|_2^2
    \quad \text{subject to } y_i(\beta^T x_i + b) \ge 1,\ \forall i.
\]

%%%%%%%%%%%%%%%%%%%%%%%%%%%%%%%%%%%%%%%%%%%%%%%%%%%%%%%%%%%%
\subsection{Soft-Margin SVM}

When data is not linearly separable, introduce slack variables $\xi_i \ge 0$:
\[
    y_i(\beta^T x_i + b) \ge 1 - \xi_i.
\]

Objective:
\[
    \min_{\beta,b,\xi} \ \frac{1}{2}\|\beta\|_2^2 + C\sum_{i=1}^N \xi_i
\]
subject to:
\[
    y_i(\beta^T x_i + b) \ge 1 - \xi_i,\qquad \xi_i \ge 0.
\]

$C > 0$ balances margin size vs.\ misclassification penalty.

%%%%%%%%%%%%%%%%%%%%%%%%%%%%%%%%%%%%%%%%%%%%%%%%%%%%%%%%%%%%
\subsection{Dual Formulation and Support Vectors}

For the soft-margin SVM, the dual problem is:

\[
    \max_{\alpha} \ Q(\alpha) =
    \sum_{i=1}^N \alpha_i
    - \frac{1}{2}\sum_{i=1}^N\sum_{j=1}^N \alpha_i \alpha_j y_i y_j (x_i^T x_j)
\]
subject to:
\[
    \sum_{i=1}^N \alpha_i y_i = 0,\quad 0 \le \alpha_i \le C.
\]

The optimal weight vector is:
\[
    \beta^* = \sum_{i=1}^N \alpha_i y_i x_i.
\]

\textbf{Support vectors} are the points with $\alpha_i > 0$; the solution depends only on them.

The decision function:
\[
    f(x) = \text{sign}\left( \sum_{i=1}^N \alpha_i y_i x_i^T x + b \right).
\]

%%%%%%%%%%%%%%%%%%%%%%%%%%%%%%%%%%%%%%%%%%%%%%%%%%%%%%%%%%%%
\subsection{Nonlinear SVM and the Kernel Trick}

We map inputs into a (possibly high-dimensional) feature space via $\phi(x)$ and apply a linear SVM in that space:
\[
    f(x) = \text{sign}\big(\beta^T \phi(x) + b\big).
\]

Instead of working explicitly with $\phi(x)$, define a kernel:
\[
    K(x_i, x_j) = \phi(x_i)^T \phi(x_j).
\]

Dual problem:
\[
    \max_{\alpha} \ Q(\alpha) =
    \sum_{i=1}^N \alpha_i
    - \frac{1}{2}\sum_{i=1}^N\sum_{j=1}^N \alpha_i \alpha_j y_i y_j K(x_i, x_j),
\]
subject to:
\[
    \sum_{i=1}^N \alpha_i y_i = 0,\quad 0 \le \alpha_i \le C.
\]

Decision function:
\[
    f(x) = \text{sign}\left( \sum_{i=1}^N \alpha_i y_i K(x_i, x) + b \right).
\]

\subsubsection*{Common Kernels}

\begin{itemize}
    \item Linear: $K(x,x') = x^T x'$.
    \item Polynomial: $K(x,x') = (1 + x^T x')^p$.
    \item Gaussian RBF: $K(x,x') = \exp\left(-\frac{\|x - x'\|^2}{2\sigma^2}\right)$.
\end{itemize}

%%%%%%%%%%%%%%%%%%%%%%%%%%%%%%%%%%%%%%%%%%%%%%%%%%%%%%%%%%%%
\section{Random Forests}
%%%%%%%%%%%%%%%%%%%%%%%%%%%%%%%%%%%%%%%%%%%%%%%%%%%%%%%%%%%%

Random forests are ensembles of decision trees.

%%%%%%%%%%%%%%%%%%%%%%%%%%%%%%%%%%%%%%%%%%%%%%%%%%%%%%%%%%%%
\subsection{Decision Trees (CART)}

Decision trees recursively partition the feature space into rectangular regions.

For regression:
\begin{itemize}
    \item At each node, choose a variable $j$ and split point $s$.
    \item Define regions:
    \[
        R_1(j,s) = \{x : x_j < s\},\quad
        R_2(j,s) = \{x : x_j \ge s\}.
    \]
    \item Choose $(j,s)$ to minimize sum of squared errors:
    \[
        \sum_{x_i \in R_1(j,s)} (y_i - \bar{y}_{R_1})^2
        + \sum_{x_i \in R_2(j,s)} (y_i - \bar{y}_{R_2})^2.
    \]
\end{itemize}

For classification:
\begin{itemize}
    \item Use impurity measures such as Gini index or deviance at each node.
    \item Choose split that most reduces impurity.
\end{itemize}

Greedy, top-down procedure; pruning can be used to avoid overfitting.

%%%%%%%%%%%%%%%%%%%%%%%%%%%%%%%%%%%%%%%%%%%%%%%%%%%%%%%%%%%%
\subsection{Random Forest Algorithm}

Random forest for classification or regression:

\begin{enumerate}
    \item For $b = 1$ to $B$ (number of trees):
    \begin{enumerate}
        \item Draw a bootstrap sample (with replacement) from the training data.
        \item Grow a tree $T_b$:
        \begin{itemize}
            \item At each node, randomly select $m$ predictors out of $p$.
            \item Find the best split among those $m$ predictors.
            \item Continue splitting until some stopping criterion (e.g.\ minimum node size) is reached; no pruning.
        \end{itemize}
    \end{enumerate}
    \item For prediction:
    \begin{itemize}
        \item Regression: average the predictions from all trees.
        \item Classification: majority vote over trees.
    \end{itemize}
\end{enumerate}

Key ideas:
\begin{itemize}
    \item \textbf{Bagging} reduces variance by averaging over bootstrap samples.
    \item Random feature selection at each split reduces correlation among trees.
\end{itemize}

%%%%%%%%%%%%%%%%%%%%%%%%%%%%%%%%%%%%%%%%%%%%%%%%%%%%%%%%%%%%
\subsection{Out-of-Bag (OOB) Error}

Because each tree trains on a bootstrap sample, about $1/3$ of the observations are \emph{out-of-bag} for that tree.

\begin{itemize}
    \item For each observation, aggregate predictions from all trees where it was OOB.
    \item Compare aggregated OOB prediction to true label.
    \item OOB error is an unbiased estimate of test error.
\end{itemize}

%%%%%%%%%%%%%%%%%%%%%%%%%%%%%%%%%%%%%%%%%%%%%%%%%%%%%%%%%%%%
\subsection{Variable Importance in Random Forests}

Permutation importance:
\begin{enumerate}
    \item Compute baseline OOB error using the original data.
    \item For each predictor $x_p$:
    \begin{enumerate}
        \item Randomly permute $x_p$ among the training samples to create $D'$.
        \item Recompute OOB error using $D'$.
        \item The increase in OOB error is the importance of $x_p$.
    \end{enumerate}
\end{enumerate}

%%%%%%%%%%%%%%%%%%%%%%%%%%%%%%%%%%%%%%%%%%%%%%%%%%%%%%%%%%%%
\subsection{Random Forests for Big Data}

\begin{itemize}
    \item Data can be very large and distributed across a cluster (e.g.\ HDFS).
    \item Each tree can be trained on a subset of the data, possibly on a separate node.
    \item Ensembles of many smaller trees can outperform a single tree trained on all data.
    \item Training and prediction are easily parallelizable.
\end{itemize}

\end{document}
